\documentclass[10pt,a4paper]{amsart}
\usepackage[margin=19mm]{geometry}
\usepackage{amsmath,amssymb,mathtools}
\usepackage[scr=rsfs,cal=euler]{mathalfa}
\usepackage{xfrac}
\usepackage[unicode,hidelinks,bookmarks=false]{hyperref}
\newcommand{\ie}{i.e.\ }
\newcommand{\cf}{cf.\ }
\newcommand{\resp}{respectively\ }
\newcommand{\Subsection}{\subsection}
\providecommand{\nobreakdash}{\nobreak\hbox{-}}
\setlength{\emergencystretch}{2em}
\multlinegap0pt
\usepackage{bm}
\usepackage{bbm}
\usepackage{dsfont}
\usepackage{xspace}
\usepackage{cancel}
\usepackage{mathtools}
\usepackage{amsthm}
\makeatletter
\@ifundefined{theorem}{\newtheorem{theorem}{Theorem}}{}
\@ifundefined{lemma}{\newtheorem{lemma}[theorem]{Lemma}}{}
\makeatother
\def\sm{\setminus}
\def\sub{\subset}
\title[The completeness of free Boolean topological groups]{The completeness of free Boolean topological groups}
\author{Ol'ga V. Sipacheva, Mikhail G. Tkachenko}
\date{Source: arXiv:2507.12289v1 (CC BY 4.0)}
\begin{document}
\maketitle

\noindent\emph{Source and licence.} The completeness of free Boolean topological groups. Version arXiv:2507.12289v1 (CC BY 4.0), distributed under the Creative Commons Attribution 4.0 International licence, as recorded in the arXiv licence metadata for this exact version and shown on the abstract page https://arxiv.org/abs/2507.12289v1. Full rights notice: licenses/arxiv-2507.12289-CC-BY-4.0.txt.\par\smallskip

\noindent\textit{Editorial scope.} Counted unit: the Lemma whose source label is \texttt{Le:clearer} in the article (source0.tex lines 295--314, source label \texttt{Le:clearer}), delivered together with the article's own complete original proof of it (source0.tex lines 316--376, one \texttt{proof} environment of 61 lines). No mathematics is written, repaired or re-derived: statement and proof are the article's own text line for line. Required context retained uncounted: eq:r1 (lines 267--270); eq:sum2 (lines 273--276). Every definition, display and label of those lines is retained unchanged. Omitted from this package and not redistributed: 1--266, 271--272, 277--294, 315--315, 377--943. Nothing inside a retained excerpt is omitted or altered: every retained body is delivered exactly as the article's lines are written, so no omission receipt of any kind accompanies this package. Citations: the retained excerpts carry no citation command at all, so no bibliography entry of the article is transcribed and no reference is redistributed. Reference adaptation: none was needed and none was made; every reference inside the delivered unit resolves inside it, so no unresolved reference or citation is produced. Printed slips kept literal: no printed word, symbol, hypothesis or number of the counted statement or of its proof was altered. Layout and class: the article's own class is \texttt{amsart} with packages including amsthm, euscript, epsf, amsmath, amssymb, latexsym, amsopn, amsfonts, fontenc, babel, courier, mathrsfs; the wrapper is the lane's pinned \texttt{amsart} editorial wrapper at 10pt on A4, which restates the theorem-like environments the retained text needs and adds the article's own macro definitions that the retained text needs (\texttt{\textbackslash sm}, \texttt{\textbackslash sub}), with extra wrapper packages (bm, bbm, dsfont, xspace, cancel, mathtools). The delivered numbering is the unit's own. Assets: no retained span contains a figure environment, a graphics-inclusion command, a PDF-page inclusion, a table or a TikZ picture; no image file is copied into this package, the package declares no asset dependency and no image file is redistributed with it. Licence: version arXiv:2507.12289v1 of this article is distributed under the Creative Commons Attribution 4.0 International licence, as recorded in the arXiv licence metadata for this exact version and shown on the abstract page https://arxiv.org/abs/2507.12289v1. A full rights notice is delivered at licenses/arxiv-2507.12289-CC-BY-4.0.txt. Characters: every retained excerpt is pure ASCII, so the delivered engine, XeLaTeX with its default Latin Modern text font, reports no missing character.
\begin{equation} 
\label{eq:r1} 
h= (x_1+y_1)+\cdots+(x_n+y_n), 
\end{equation} 

\begin{equation}
\label{eq:sum2}
\sum_{i=1}^n \varrho(x_i,y_i). 
\end{equation} 

\begin{lemma}\label{Le:clearer}
Suppose that $\varrho$ is a pseudometric on a set $X$ and 
$h=x_1+\cdots+x_n$ is an element of $B_a(X)\sm\{e\}$ 
with pairwise distinct $x_1,\ldots,x_n\in X\sm\{e\}$. Then 
there exists a representation
\begin{eqnarray*}
h=(u_1+v_1)+\cdots+(u_k+v_k) 
\end{eqnarray*}
such that $u_1,v_1,\ldots,u_k,v_k$ are  pairwise distinct and 
$$
\widehat{\varrho}(h,e)=\sum_{i=1}^k \varrho(u_i,v_i). 
$$
Moreover, if $n$ is even, then $2k=n$ and $\{u_1,v_1,\ldots,u_k,v_k\}=
\{x_1,\ldots,\allowbreak x_n\}$, and if $n$ is odd, then $2k=n+1$, 
$\{u_1,v_1,\ldots,u_k\}=\{x_1,\ldots, x_n\}$, and $v_k=e$. 

In particular, the quantity $N_\varrho(h)$, which is 
the infimum of the sums $(\ref{eq:sum2})$ over all possible representations 
of $h$ of the form \eqref{eq:r1}, is in fact the minimum of those sums. 
\end{lemma}

\begin{proof}
Let 
\begin{equation}
\label{eq:r3}
h = (z_1+t_1) + \cdots + (z_p+t_p)
\end{equation}
be a representation of $h$ with $z_i,t_i\in X$ for each $i\leq p$.
First, we claim that certain cancellations in $(\ref{eq:r3})$ yield another 
representation 
\begin{equation}
\label{eq:r6}
h = (a_1+b_1) + \cdots + (a_q+b_q)
\end{equation}
such that $q\leq p$, $\{a_1,b_1,\ldots,a_q,b_q\}\sub 
\{z_1,t_1,\ldots,z_p,t_p\}$, the elements $a_1,b_1,\ldots,a_q,b_q$
are pairwise distinct, and 
\begin{equation}
\label{eq:r9}
\sum_{i=1}^q \varrho(a_i,b_i)\leq 
\sum_{i=1}^p \varrho(z_i,t_i).
\end{equation}
Clearly, we may assume that $z_i\neq t_i$ for each $i\leq p$; otherwise,
we delete the summand $z_i+t_i$ from the representation $(\ref{eq:r3})$
of $h$, affecting neither $h$ nor the sum on the right-hand side of 
$(\ref{eq:r9})$. Suppose that $i\ne j$ and the summands $(z_i+t_i)$
and $(z_j+t_j)$ in $(\ref{eq:r3})$ contain equal elements. We may assume 
without loss of generality that $z_i=z_j$, because the group $B_a(X)$ is 
Abelian. We replace the two summands $(z_i+t_i)$ and $(z_j+t_j)$ in 
$(\ref{eq:r3})$ with the single summand $(t_i+t_j)$, thus obtaining a shorter 
representation of $h$. In view of the triangle inequality we have 
$\varrho(t_i,t_j)\leq \varrho(t_i,z_i) + \varrho(z_j,t_j)$ (here $z_i=z_j$), 
so that the sum of distances for the new representation of $h$ does not 
increase. Proceeding, we obtain a reduced representation of $h$ of the form 
$(\ref{eq:r6})$ satisfying condition $(\ref{eq:r9})$. In particular, all 
elements $a_1,b_1,\ldots,a_q,b_q$ of this representation are pairwise 
distinct and $\{a_1,b_1,\ldots,a_q,b_q\}\sub \{z_1,t_1,\ldots,z_p,t_p\}$, 
because the transformations of the given representation of $h$ performed above 
introduce no new elements of~$X$.

Note that $h=x_1+\dots+x_n$ is an expansion of $h$ in the basis $X\setminus 
\{e\}$. Since the elements $a_1,b_1,\ldots,a_q,b_q$ are 
pairwise distinct, it follows that at most one of them equals $e$ and 
the remaining ones belong to $X\setminus \{e\}$. 
Therefore, $\{a_1,b_1,\ldots,a_q,b_q\}\subset \{x_1,\dots, x_n\}\cup\{e\}$.  
Clearly, $h$ has only finitely many representations of the form \eqref{eq:r6} 
satisfying this condition. Since inequality \eqref{eq:r9} holds for all such 
representations, it follows from the arbitrariness of the representation 
\eqref{eq:r3}  that one of them, say, $h = (u_1+v_1)+\dots+ (u_k+v_k)$, 
satisfies
$$ 
\widehat{\varrho}(h,e) =N_\varrho(h)=\sum_{i=1}^q \varrho(a_i,b_i). 
$$


Since $u_1,v_1,\ldots,u_k,v_k$ are pairwise distinct
and $\{u_1,v_1,\ldots,u_k,v_k\}\sub\{x_1,\ldots,x_n,e\}$, we have $n\leq 2k\leq 
n+1$. Therefore, $n=2k$ if $n$ is even and $n+1=2k$ if $n$ is odd. In the 
latter case, exactly one of the elements 
$u_1,v_1,\ldots,u_k,v_k$ equals $e$, and we can assume that this is $v_k$, 
because $B_a(X)$ is Abelian. This completes the proof of the lemma. 
\end{proof}

\end{document}
